\chapter{Conclusion}\label{ch:conclusion}
Digital Mohalla addresses one question: whether a proposed neighbourhood loan has enough trustworthy evidence, when that evidence is sparse, role-dependent, indirect and open to manipulation. Research on reputation systems showed that averages hide uncertainty and that newcomers need a policy, not a score; research on collusion showed that a signal must separate rings from ordinary two-sided reviewing. These findings became the requirements and the design: four evidence states instead of a trust number, a recency-weighted Beta posterior with an interval, capped indirect evidence, corroborated manipulation signals, and a human moderator for every hold and dispute.

Iteration 1 built that core and the lifecycle that feeds it. The lifecycle is event-sourced, replayable and protected by the database against invalid transitions and double bookings; the engine runs on real outcomes and explains every result; disputes and holds reach a moderator whose decision becomes evidence. 623 backend and 29 frontend tests pass, all seven scenarios required at mid-term are correct in state and explanation, and the system is deployed from \code{main} after every successful CI run. Generative AI assistance is declared per module, with the team's own design decisions and reviews recorded alongside it. The main risk is that the engine's parameters are still hypotheses. Iteration 2 completes the core modules and builds the simulator and baselines, so that FYP-2 can test, with numbers, whether the engine's uncertainty leads to better lending decisions than the averages it replaces.
